% THETA3-B27 contribution by /root/definite_bridge, 21 November 2026.
\subsection{THETA3-B27: the actual mixed source has multiplicity-free Hecke sectors}\label{B27}

For the actual rank-two source of B26, the Hermitian multiplicity
matrices in (B26r) can be reduced further: every multiplicity is one.
Consequently the Gram of the original fixed kernel depends on finitely
many positive scalar periods, one for each source Hecke sector.
This assertion does not require the source motive to be CM and does
not assert that only one sector occurs.  The argument uses only the
generic pure-inner-form case of the unitary multiplicity formula,
after excluding the nongeneric rank-two parameter directly.

\paragraph{The fixed archimedean representation.}
Retain $S=\U(P)$, the nonempty switch set $I$, the complex orientations,
the full central character, and the algebraic line $\mathscr L$ of B26.
At each $\tau\in I$ the source representation generated by a nonzero
holomorphic coefficient vector restricts to the holomorphic
discrete series of $\SU(1,1)$ of lowest compact weight three.
At every other real place its complete compact type is the fixed
one-dimensional determinant character.  In particular the entire
archimedean representation, denoted $\sigma_\infty$, is fixed, and
its specified holomorphic compact type has dimension one.

Here this conclusion includes the lowering-operator condition,
not merely a coincidence of compact weights.  The actual Fock
vacuum in B26 is killed by the lowering root.  A unitary
lowest-weight $\mathfrak{sl}_2$ module of lowest weight three is the
discrete series of that weight: successive raising vectors have
weights $3+2j$, and the commutator relations and positivity determine
the module and its invariant inner product up to a scalar.  Thus
the holomorphic sections in question cannot contribute a different
unitary archimedean representation with the same isolated compact
weight.  The fixed central character determines the extension
from $\SU(1,1)$ to $\U(1,1)$, since
$\U(1,1)=(\SU(1,1)\times\U(1))/\{\pm1\}$.
At a compact place the compact group is the entire real group,
so its specified scalar type already determines the representation.
The simultaneous contragredient convention replaces all these
holomorphic choices by their prescribed duals and leaves the
argument unchanged.

Because $P$ is definite at $\mu$, it is anisotropic over $F$.
Its unitary group and its center have compact adelic quotients.
Thus the source functions belong to the discrete automorphic
spectrum, with no continuous-spectrum issue.

\paragraph{A derived-invariance lemma at one split finite place.}
Let $D=\SU(P)$.  This is a simply connected absolutely almost
simple $F$-group of type $A_1$.  Let $v$ be a finite place split
in $E$, so $D(F_v)\simeq\mathrm{SL}_2(F_v)$ is noncompact.
Strong approximation gives
\[
       \overline{D(F)}=D(\A_F^{\{v\}}).
 \tag{B27a}
\]
The notation on the right includes every archimedean place.
This is precisely Platonov--Rapinchuk, \emph{Algebraic Groups
and Number Theory}, Theorem 7.12, printed p.427; the final part
of its proof on pp.431--432 removes the condition that the set
of omitted places contain the archimedean ones.  Their theorem
therefore applies with the singleton set $\{v\}$.

Suppose a continuous automorphic function $f$ on $S(F)\backslash
S(\A_F)$ is right invariant under $D(F_v)$.  Then it is right
invariant under all of $D(\A_F)$.  To prove this at an arbitrary
base point $g=(g^v,g_v)$, take $\gamma\in D(F)$.  Normality of
$D(F_v)$ in $S(F_v)$ and the assumed right invariance give
\[
 f(\gamma^v g^v,g_v)
   =f(\gamma^v g^v,\gamma_v g_v)=f(g).
 \tag{B27b}
\]
For any $d^v\in D(\A_F^{\{v\}})$, choose such rational
$\gamma$ converging to $d^v$ away from $v$ and apply continuity.
This proves left $D(\A_F^{\{v\}})$ invariance.  At $v$ the
assumed right invariance is also left invariance by normality.
Thus $f$ is left invariant under $D(\A_F)$, and normality once
again makes this equivalent to right invariance.  No condition
on the possibly unbounded $v$-components of the approximating
rational elements is needed.

\paragraph{Exclusion of the only nongeneric rank-two parameter.}
Use the formal global parameters and near-equivalence decomposition
of Kaletha--M\'inguez--Shin--White, \emph{Endoscopic classification
of representations: inner forms of unitary groups I},
Section 1.3 and Corollary 3.3.2, printed pp.63--70 and 151.
The latter decomposes the discrete spectrum according to
parameters using their unramified Satake classes; it is available
before the full local classification of nongeneric inner-form
packets.  In rank two the only possible nongeneric parameter is
\[
                         \psi=\chi\boxtimes S_2,
 \tag{B27c}
\]
where $\chi$ is a conjugate-self-dual automorphic character of
$\A_E^\times/E^\times$ and $S_2$ is the two-dimensional
irreducible representation of the Arthur $\mathrm{SU}(2)$.
Indeed the parameter is a sum of blocks of dimensions
$m_i n_i$; a nontrivial Arthur factor has $n_i\ge2$ and
already exhausts the rank two.  Parameters which are sums
of rank-one blocks with trivial Arthur factors are generic,
even when the sum is not discrete.

Choose a finite place $v$ split in $E$ at which the automorphic
representation, the parameter, and the fixed base-change splitting
character are unramified.  Such places exist outside any finite
set.  Here $S(F_v)\simeq\mathrm{GL}_2(F_v)$, and (B27c)
has normalized Satake eigenvalues
\[
              \alpha_v q_v^{1/2},\qquad
              \alpha_v q_v^{-1/2}
 \tag{B27d}
\]
for a nonzero scalar $\alpha_v$.  All fixed splitting-character
factors multiply both eigenvalues by the same scalar.  KMSW's
formula (1.3.2), printed p.64, and the localization in
Section 1.3.3 give exactly these exponents.  The unique spherical
irreducible representation with these eigenvalues is an
unramified character composed with the determinant.  This can
also be checked without a nonsplit packet theorem: the trivial
representation of $\mathrm{GL}_2(F_v)$ is the spherical
constituent of normalized induction of
$|\cdot|^{1/2}\otimes|\cdot|^{-1/2}$, and a determinant twist
multiplies both Satake eigenvalues by $\alpha_v$.

Consequently the local representation is trivial on $D(F_v)$.
The lemma (B27a)--(B27b) makes its automorphic realization
trivial on $D(\A_F)$.  This contradicts the lowest weight three
at any place in $I$.  Hence no actual source constituent has
parameter (B27c): every constituent is generic.

This exclusion uses a split finite place and strong approximation.
It does not infer genericity merely from an unproved compatibility
of the archimedean infinitesimal character with a nongeneric
inner-form Arthur packet.

\paragraph{The precise multiplicity formula that now applies.}
The group $S$ is a pure inner form of the quasi-split unitary
group: it is defined by a Hermitian space over $E$, rather than
by a Hermitian space over a nontrivial division algebra.
KMSW Section 0.3.3, printed p.18, proves exactly this criterion
by the relevant Galois cohomology sequence and the Brauer class.
The generic pure-inner-form case of their Theorem 5.0.5,
printed pp.205--206, is therefore applicable to every actual
source constituent.  Its two local hypotheses are expressly
proved for generic parameters and pure inner twists in the
preceding paragraph.  The theorem gives zero contribution for
a nondiscrete generic parameter and, for a discrete generic
parameter, a direct sum in which each permitted automorphic
representation occurs once.

We use this proved case, not the unrestricted starred statement
of Theorem 1.7.1.  The scope paragraph after that starred
statement, printed p.89, explicitly distinguishes the proved
generic pure-inner case from the other cases reserved for
sequels.  We also make no assertion that the local components
at every finite place are tempered; their generic global
parameter suffices here.

It follows that every irreducible automorphic representation
with the actual infinity type $\sigma_\infty$ occurs with
multiplicity one.  At a fixed finite compact open $K_S$, the
full finite-level Hecke module of holomorphic sections is
therefore multiplicity free as well:
\[
          H^0(Y,\mathscr L)
              =\bigoplus_{\sigma}
                    \sigma_f^{K_S},
 \tag{B27e}
\]
where zero summands are omitted and the equality refers to the
fixed full adelic realization and central character of B26.
Every summand on the right is a simple module for the full
$K_S$-bi-invariant Hecke algebra.  Two different representations
on the right cannot yield isomorphic nonzero Hecke modules.

For completeness, these finite-Hecke assertions are elementary
idempotent facts.  If $e_{K_S}$ is normalized compact averaging
and $V$ is an irreducible smooth representation with
$e_{K_S}V\ne0$, a nonzero invariant subspace of $e_{K_S}V$
generates $V$ under the group, and averaging that generated
space back into $e_{K_S}V$ proves its simplicity.  Conversely
a simple $e_{K_S}\mathcal H e_{K_S}$-module determines the unique
simple quotient generated by it: in the induced module, quotient
by the largest submodule annihilated by all operators
$e_{K_S}\mathcal H$; any further nonzero submodule contains the
generating $e_{K_S}$-module.  Thus the finite representation is
determined by its nonzero fixed-vector Hecke module.  Finally,
two global representations with the same finite representation
and the already fixed $\sigma_\infty$ are the same restricted
tensor product.  The selected archimedean type contributes
one dimension, not another multiplicity.

\paragraph{The original fixed-kernel Gram now has only scalar sectors.}
Let $\mathcal F_{\overline\Q}$ be the actual algebraic Hecke
saturation constructed in B26.  It is a submodule of (B27e),
so every $m_\lambda$ in (B26p) is one.  In the algebraic
decomposition and with the same positive algebraic matrices
$G_\lambda$ as in (B26q), there are actual positive real
numbers $c_\lambda$ such that
\[
 \begin{aligned}
  \mathcal F_{\overline\Q}&=\bigoplus_\lambda V_\lambda,
       &G_S&=\bigoplus_\lambda c_\lambda G_\lambda,\\
  \widehat W&=|\gamma|^2
       \sum_\lambda c_\lambda C_\lambda^*G_\lambda C_\lambda,
       &\gamma&=\frac{C_x}{t_x t_y}.
 \end{aligned}
 \tag{B27f}
\]
Here $C_\lambda$ is the algebraic coefficient matrix of
the normalized actual source coefficients in $V_\lambda$.
Every Haar, theta-group, tangent, and constant CM-line
factor in $\gamma$ and the source periods is retained from
B26.  This is the Gram of the original kernel.  Its remaining
freedom is at most one positive scalar per sector, rather
than an arbitrary positive Hermitian matrix on each
multiplicity space.

The central idempotents $p_\lambda$ of the finite algebraic
Hecke image are algebraic and are orthogonal projections for
the actual Petersson form.  Each can be represented by a
finite algebraic linear combination of finite double-coset
operators.  Apply that same operator to the source variable
of the original kernel:
\[
 K_\lambda(g,h)=p_\lambda^h K(g,h),\qquad
 K=\sum_\lambda K_\lambda.
 \tag{B27g}
\]
These are genuine theta kernels on the same $P,V$ with
the same archimedean cochain.  Indeed
$K_{\omega(s)\Phi}(g,h)=K_\Phi(g,hs)$ for a finite source
element $s$, and convolution by a finite double-coset
operator is the corresponding finite algebraic combination
of finite Weil translates.  Its coefficients are algebraic
in the finite rational Weil structure used in B25, with
the same splitting-character values retained.  No direct
sum of unrelated Hermitian spaces is introduced.

Let $T_\lambda$ be the bilinear theta integral with kernel
$K_\lambda$.  Orthogonality in the source, with exactly
the Hilbert-space convention of B17, gives
\[
 T_\lambda T_\kappa^*=0\quad(\lambda\ne\kappa),
       \qquad TT^*=\sum_\lambda T_\lambda T_\lambda^*.
 \tag{B27h}
\]
The first identity is the integral of a source coefficient
in one sector against the conjugate of a coefficient in
another.  In particular it is valid for the bilinear
theta-integral convention; no unnoticed inversion of a
source character is being used.

There is also an explicit matrix form.  Use B25's algebraic
target basis, with target Gram
$G_T=\nu_G|P_\chi|^2G_B$, and put
\[
 D_\lambda=(C_\lambda^*G_\lambda C_\lambda)^{\mathsf t},
       \qquad A_\lambda=D_\lambda G_B.
 \tag{B27i}
\]
The transpose is the one required by B17's
$W_{ab}=\int e_a\overline{e_b}$.  Thus
\[
 T_\lambda T_\lambda^*
   =\nu_G|P_\chi|^2|\gamma|^2c_\lambda A_\lambda.
 \tag{B27j}
\]
The matrix $A_\lambda$ is algebraic and is positive
semidefinite and self-adjoint for $G_B$.  Its image
$E_\lambda$ is algebraic.  On that image it is positive
definite, so its inverse is algebraic.  Equivalently,
extend that inverse by zero on the $G_B$-orthogonal
complement.  B23 realizes this algebraic matrix by a
correspondence from the finite CM source.  Therefore the
actual B17 inverse in a single nonzero sector is
\[
 (T_\lambda T_\lambda^*|_{E_\lambda})^{-1}
   =\frac{1}{\nu_G|P_\chi|^2|\gamma|^2c_\lambda}
                          (A_\lambda|_{E_\lambda})^{-1}.
 \tag{B27k}
\]
This closes its matrix-valued algebraicity issue within
each actual sector.  The one remaining scalar is its
Petersson period with the already specified normalization.

For an original nonzero lift $Tu$, the decomposition
$Tu=\sum_\lambda T_\lambda u$ has at least one nonzero
summand.  There is no assertion that the original lift
itself lies in only one $E_\lambda$, or that the target
images for different sectors are orthogonal or disjoint.
The original full $TT^*$ remains the weighted sum in
(B27h)--(B27j).  Its inverse is projectively algebraic
if the relevant $c_\lambda$ have algebraic ratios;
that comparison has not been proved here.  In particular
this result does not identify a single prescribed
\texttt{MainX} tuple or supply its arithmetic constants.

\paragraph{Primary-source and proof scope.}
The actual source weights, line, and common coherent
scalar are B26, with the exact 032 and local-theta
sources recorded there.  For the new global input the
whole cited KMSW statements and their immediate proofs
were read: Section 0.3.3, Section 1.3.2--1.3.4, Corollary
3.3.2 and its decomposition, the scope paragraph after
Theorem 1.7.1, and all of Theorem 5.0.5 and its proof.
The strong approximation source is the original authors'
Theorem 7.12 with its complete proof, including the
singleton finite-place scope.  Its application (B27b),
the rank-two exclusion, the fixed-infinity Hecke
multiplicity argument, and the displayed kernel and
matrix consequences were proved above.  No unproved
nongeneric inner-form packet theorem is needed.
